% 彗星（综述）
% license CCBYSA3
% type Wiki

本文根据 CC-BY-SA 协议转载翻译自维基百科\href{https://en.wikipedia.org/wiki/Comet}{相关文章}。

\begin{figure}[ht]
\centering
\includegraphics[width=8cm]{./figures/45ca94c34375617b.png}
\caption{} \label{fig_Comet_1}
\end{figure}
彗星是一类由冰组成的、体积较小的太阳系天体或星际天体。当其经过接近太阳的区域时，会因受热而开始释放气体，这一过程称为释气作用（outgassing）。这一过程会在彗核周围产生一个延展的、未被引力束缚的大气层，即彗发（coma），并有时形成由彗发中吹离的气体和尘埃构成的彗尾。这些现象源于太阳辐射以及向外流动的太阳风等离子体对彗核的作用。彗核的大小从数百米到数十公里不等，由松散聚集的冰、尘埃和小型岩石颗粒组成。彗发的尺度可达地球直径的 15 倍，而彗尾的长度甚至可能超过一个天文单位。如果距离足够近且亮度足够高，人类在无望远镜的情况下即可从地球上观测到彗星，其在天空中的角距可达 30°（相当于 60 个月亮）。自古以来，彗星已被众多文化和宗教观察并记录。

彗星通常具有高离心率的椭圆轨道，其轨道周期范围极广，从数年到可能长达数百万年。短周期彗星来源于海王星轨道之外的柯伊伯带或相关的散射盘。长周期彗星则被认为起源于奥尔特云，一个由冰质天体组成的球形云层，其范围从柯伊伯带外侧延伸至通往最近恒星方向一半的距离[2]。长周期彗星通常因经过恒星的引力扰动或银河潮汐效应而被触发向太阳方向运动。双曲轨道彗星可能只会一次穿越内太阳系，随后被抛入星际空间。彗星的出现称为一次显现（apparition）。

多次接近太阳的灭绝彗星会失去几乎全部挥发性冰与尘埃，并可能逐渐呈现出类似小型小行星的外观[3]。一般认为，小行星与彗星的起源不同：小行星形成于木星轨道以内，而非太阳系外部[4][5]。然而，主带彗星（main-belt comets）以及活跃的半人马族小天体的发现，使小行星与彗星之间的界限变得模糊。21 世纪初，一些具备长周期彗星轨道却呈现内太阳系小行星特征的小天体被称为曼岛彗星（Manx comets）。它们仍被归类为彗星，例如 C/2014 S3（PANSTARRS）[6]。在 2013 年至 2017 年之间，共发现了 27 颗曼岛彗星[7]。

截至 2021 年 11 月，共已知 4,584 颗彗星[8]。然而，这仅占可能存在的彗星总量的一小部分，因为太阳系外部（奥尔特云）中类彗星天体的储库估计约为一万亿个[9][10]。平均每年约有一颗彗星可被肉眼看见，尽管其中许多较为暗淡、外观并不特别醒目[11]。特别明亮的彗星被称为“伟大彗星”（great comets）。人类已利用无人探测器访问过彗星，例如美国 NASA 的“深度撞击号”（Deep Impact），其通过撞击坦普尔 1 号彗星形成陨坑以研究其内部结构；又如欧洲航天局的“罗塞塔号”（Rosetta），它成为首个在彗星上着陆的机器人探测器[12]。
\subsection{词源}
\begin{figure}[ht]
\centering
\includegraphics[width=8cm]{./figures/6a3173d4cf8911ea.png}
\caption{公元 729 年的《盎格鲁－撒克逊编年史》以及尊敬的比德（the Venerable Bede）的著作中均提到了彗星。} \label{fig_Comet_2}
\end{figure}
单词“comet（彗星）”源自古英语 cometa，其来源为拉丁语 comēta 或 comētēs，而这些拉丁词又是对希腊语 κομήτης 的罗马化形式，意为“留着长发的（事物）”。《牛津英语词典》指出，术语（ἀστὴρ） κομήτης 在希腊语中早已表示“长发之星，即彗星”。κομήτης 源自动词 κομᾶν（koman），意为“留长发”，而该动词又派生自名词 κόμη（komē），意为“头发”，并被用来表示“彗尾”[15][16]。

彗星的天文学符号（在 Unicode 中表示）为 U+2604 ☄ COMET，由一个小圆盘及其后三条类似发丝的延伸部分组成[17]。
\subsection{物理特征}
\begin{figure}[ht]
\centering
\includegraphics[width=8cm]{./figures/37fee7dd52548b31.png}
\caption{彗星的结构} \label{fig_Comet_3}
\end{figure}
\subsubsection{彗核}
\begin{figure}[ht]
\centering
\includegraphics[width=8cm]{./figures/43510b130873771c.png}
\caption{在航天器飞越期间拍摄的 103P/Hartley 彗核图像。该彗核长度约为 2 千米。} \label{fig_Comet_4}
\end{figure}
彗星的固体核心结构被称为彗核。彗核由岩石、尘埃、水冰以及冻结的二氧化碳、一氧化碳、甲烷和氨等物质的混合体构成[18]。因此，根据弗雷德·惠普尔（Fred Whipple）的模型，它们常被形象地称为“肮脏的雪球”（dirty snowballs）[19]。尘埃含量更高的彗星则被称为“冰质尘球”（icy dirtballs）[20]。“冰质尘球”这一术语源于 2005 年 7 月 NASA“深度撞击号”（Deep Impact）任务向坦普尔 1 号彗星（Comet 9P/Tempel 1）发送“撞击器”探测器并观测其碰撞后的结果。2014 年的一项研究表明，彗星类似于“炸冰淇淋”，即其表面由密集的晶体冰与有机化合物混合构成，而内部冰层则更冷且密度更低[21]。

彗核表面通常干燥、多尘或岩质，表明冰层隐藏在厚度数米的表层壳体之下。彗核中含有多种有机化合物，其中可能包括甲醇、氰化氢、甲醛、乙醇、乙烷，以及可能更复杂的分子，如长链烃类与氨基酸[22][23]。2009 年确认，NASA “星尘号”（Stardust）任务回收的彗星尘埃中发现了氨基酸甘氨酸[24]。2011 年 8 月，一份基于 NASA 对地球陨石研究的报告提出，DNA 与 RNA 的组成单元（腺嘌呤、鸟嘌呤及相关有机分子）可能形成于小行星和彗星[25][26]。

彗核外层表面具有极低反照率，使其成为太阳系中反射率最低的天体之一。贾托号（Giotto）探测器发现，哈雷彗星（1P/Halley）的彗核仅反射约四个百分点的入射光[27]；而“深空一号”（Deep Space 1）发现波罗雷彗星（Comet Borrelly）的表面反射率甚至不足 3.0\%[27]。相比之下，沥青可反射约 7\%。彗核表面的暗色物质可能由复杂的有机化合物构成。太阳加热会驱散较轻、易挥发的化合物，而留下更大的有机分子，这些物质通常极为黝黑，如焦油或原油。彗核表面的低反照率使其能够吸收驱动释气过程的热量[28]。

已观测到的彗核半径可达 30 千米（19 英里）[29]，但要准确测定彗核的大小并不容易[30]。322P/SOHO 的彗核直径可能只有 100–200 米（330–660 英尺）[31]。尽管观测仪器的灵敏度不断提高，但仍缺乏对更小彗星的探测，这使部分研究者推测直径小于 100 米（330 英尺）的彗星数量可能确实稀少[32]。已知彗星的平均密度估计约为 $0.6 g/cm^3$（0.35 oz/cu in）[33]。由于质量较低，彗核不会在自引力作用下成为球形，因此呈现不规则形状[34]。
\begin{figure}[ht]
\centering
\includegraphics[width=8cm]{./figures/7e9645911c29b572.png}
\caption{彗星 81P/Wild 在光照面和阴暗面都呈现出喷流，其表面起伏鲜明，且十分干燥。} \label{fig_Comet_5}
\end{figure}
大约六个百分点的近地小行星被认为是不再发生释气作用的彗星灭绝彗核[35]，其中包括 14827 Hypnos 与 3552 Don Quixote。

“罗塞塔号”（Rosetta）与“菲莱”（Philae）着陆器的探测结果显示，67P/丘留莫夫－格拉西缅科彗星（67P/Churyumov–Gerasimenko）的彗核不存在磁场[36][37]，这表明磁性可能并未在早期微行星形成过程中发挥作用。此外，罗塞塔号上的 ALICE 光谱仪确定，在距离彗核约 1 千米范围内，由太阳辐射对水分子进行光电离所产生的电子，而非此前认为的来自太阳的光子，是导致彗核释放出的水与二氧化碳分子在彗发中发生分解的主要因素[38][39]。菲莱着陆器上的仪器在彗星表面检测到至少十六种有机化合物，其中有四种（乙酰胺、丙酮、异氰酸甲酯和丙醛）是首次在彗星上被发现的[40][41][42]。
\begin{figure}[ht]
\centering
\includegraphics[width=10cm]{./figures/8c2117a476d375ac.png}
\caption{} \label{fig_Comet_6}
\end{figure}
